\section{Normative gap register}
\label{sec:gaps}

\begingroup\scriptsize
\linespread{0.96}\selectfont

A gap is not permission to infer a rule.  Unless the fixed proof system supplies
the closure named below, the verifier \must{} reject the affected instance.
Rows carrying a closure qualifier distinguish a fixed core rule from
remaining subprotocol, implementation, fixture, or provenance work.

\renewcommand{\arraystretch}{1.02}
\setlength{\tabcolsep}{3pt}
\setlength{\LTleft}{0pt}
\setlength{\LTright}{\fill}
\begin{longtable}{@{}p{0.52in}
  >{\raggedright\arraybackslash}p{2.22in}
  >{\raggedright\arraybackslash}p{2.22in}
  >{\raggedright\arraybackslash}p{2.22in}@{}}
\toprule
\textbf{ID}&\textbf{Source omission or ambiguity}&
\textbf{Behavior of this specification}&\textbf{Required closure}\\
\midrule
\endfirsthead
\toprule
\textbf{ID}&\textbf{Source omission or ambiguity}&
\textbf{Behavior of this specification}&\textbf{Required closure}\\
\midrule
\endhead
\textbf{G1}&
The GKR reduction's phase-\(c\) claimed sum \(S_{c,k}\) is
transmitted, but the core reduction does not by itself prescribe its
transcript treatment.  If
\(\eq(r_{x,k},z_{x,k})=0\), the phase link does not determine
that sum.  The current Rust merged-forest verifier uses the claimed sum only
after its phase-\(c\) challenges and does not absorb it.&
The PCS-level schedule treats GKR atomically.  The fixed GKR protocol must
absorb the domain-code \code{3001} payload with the encoding fixed by
Section~\ref{sec:wire-events} for \(S_{c,k}\) after the phase-link
check and before the first phase-\(c\) round message.  Its already-known
\(S_{x,k}=e_k\) field is either omitted or carried as a
checked, non-absorbed hint.&
An interoperable GKR protocol must pin its remaining events in the reserved
\code{3002--31ff} ranges, and the current verifier must insert the
domain code \code{3001} absorb at the stated deadline; alternatively the protocol must
specify a canonical degenerate-branch rejection rule.\\
\addlinespace
\textbf{G2}&
The paper-level edges are
\(\Pi_{\rm GKR}:\mathcal R_{\rm PESAT}[\bm h]
\to\mathcal R_{\rm sum}^{h}\),
\(\Pi_{\rm MLE}:\mathcal R_{\rm sum}^{h}
\to\mathcal R_{\rm eval}^{h}\), and
\(\Pi_{\mathsf{iopp}}:\LIN(\K,\mathcal C,\Pack)
\to\mathcal R_{\rm triv}\).
The concrete path exposes the GKR sum claim before the MLE claim.&
The proof system delegates \(\Pi_{\rm GKR}\) and \(\Pi_{\rm MLE}\) to the fixed
GKR and row-LinCheck ledgers.  Ring switching and recursive opening are the
internal expansion of
\(\Pi_{\mathsf{iopp}}\).&
A proof system claiming the
abstract theorem must fix coordinate order and prove both that its concrete
\(\Pi_{\rm GKR};\Pi_{\rm MLE}\) realizes the named MLE successor and that its internal
opening path realizes \(\Pi_{\mathsf{iopp}}\) from that LIN relation to
\(\mathcal R_{\rm triv}\).\\
\pagebreak[4]
\addlinespace
\textbf{G3}&
A nontrivial virtualization matrix may yield the general coefficient
\(\bm M_{\K}^{\mathsf T}\bm a\), so an MLE-only ring switch would not cover
the virtualized relation.&
The verifier applies the public adjoint without a message, reshapes its
coefficient at the $128$-bit packing boundary, and derives the canonical
decomposition
$A_f[v,j]=\sum_\theta A_\theta[v]B_\theta[j]$.
The structured ring switch consumes that decomposition directly; the
canonical $128$-term fallback covers every coefficient, while configured
translated-equality decompositions may use fewer terms.&
When virtualization is enabled, the implementation must type $\bm M$ and
compute the public transpose action and ring-switch decomposition in the fixed
coordinate order.  It must verify the decomposition identity and provide the
configured evaluators for every $B_\theta$ and the final packed coefficient.
No post-adjoint coefficient LinCheck is permitted.
\\
\addlinespace
\textbf{G4}\\[-2pt]\textsc{rule fixed}&
An arbitrary or zero input functional might appear to make the derived row
factor vanish.&
The input Sumcheck derives $a_r=\eq_{\Fq}(r,\bm s_r)$ and places $\kappa$ in
the column factor $a'_c=\kappa\eq_{\Fq}(c,\bm s_c)$.  Since
$\sum_ra_r=1$, the row vector is nonzero even when $\kappa=0$.&
Implementations and fixtures must derive the factors in exactly this order,
including a $\kappa=0$ case; moving $\kappa$ into the row factor is
nonconforming.\\
\addlinespace
\textbf{G5}\\[-2pt]\textsc{rule fixed}&
A PIOP may produce an arbitrary coefficient on $\bm h$, while the
integer-fold core accepts only a separable row/column coefficient on $H$.&
F2Z accepts arbitrary $(\bm v,\mu)$ and runs the fixed
input-Sumcheck phase of $\Pi_1$.  It derives
$\kappa=\widetilde{\bm v}(\bm s_h)$ and the new equality weights
$W[r,c]=\kappa\eq_{\Fq}((r,c),\bm s_h)=a_ra'_c$; it never factorizes the
original $V$.&
Implement domain codes \code{1002/1101/1003}, the unbiased $\Fq$ sampler, terminal
coefficient evaluator, $2m_h/|\Fq|$ error term, and positive and negative
transcript fixtures before claiming conformance.\\
\addlinespace
\textbf{G6}\\[-2pt]\textsc{wire fixed}&
Semantic transcript encoding alone does not choose interoperable bytes, the outer
two-stream container, a recursive-opening configuration, or a proof-of-work
predicate and difficulty.  The current one-stream codec writes \(\mu\) but
absorbs its derived image \(g^\mu\), which cannot be represented by the selected
two-stream wrapper's write-and-absorb NARG operation.&
Section~\ref{sec:wire-format} fixes the wire/transcript question: it gives the
primitive, vector, record, and outer encodings; numeric
domain registry; transcript transition; NARG/hint replay; recursive Ligerito
integer tables and event ledger; Merkle row and multiproof bytes; PoW context,
nonce, predicate, schedule, and difficulty; all four length metrics; and exact
nested and outer exhaustion.  It sends and absorbs
the mandatory domain code \code{1002} input-Sumcheck rounds, samples domain code \code{1101},
and absorbs the derived domain code \code{1003} seam before the product claim.  It
then sends and absorbs
\(\nu=g^{\mu}\) under domain code \code{2001}, then checks the domain code \code{2081}
\(\mu\) hint before domain code \code{2101}.  A missing or
unknown GKR event ledger or constrained-code construction is a
hard reject, not an inferred encoding.  It also fixes the ordered
domain-code \code{4001} structured ring vectors, their anchor check before
domain code \code{4101}, and the domain-code \code{4102} term-batching
challenges.  Domain codes \code{4002/4103} remain reserved.&
\textbf{Owner: John Wu.}  Port Section~\ref{sec:wire-format} to the host and
recursive verifiers, remove the opaque bincode payload, and publish the
required complete proof, transcript, Merkle, PoW, and negative fixtures.  The
current Rust codec remains nonconforming and its legacy bytes are not v1.
End-to-end interoperability may be claimed only after the fixed GKR protocol
has published its byte ledger and the fixed code specification gives a
complete canonical RS/LCH construction
(including domain, basis, padding, and coordinate order) with independent
commitment/root vectors.  The wire namespaces and PoW rules themselves are not
left to implementation inference.\\
\addlinespace
\textbf{G7}\\[-2pt]\textsc{rule fixed}&
The current encoded outer row LinCheck includes
\(\mathtt{presum.claimed\_sums}\), and current recursive/host decoders may
accept unused roots, nonces, recursive proofs, or suffix bytes.  The presum is
currently checked only after the MultiDegree challenges.&
Wire format v1 omits that redundant field and derives
\([e_d+1]\); the reserved legacy domain code \code{3081} rejects.
Sections~\ref{sec:wire-merkle}, \ref{sec:wire-events},
\ref{sec:wire-recursive-ledger}, and~\ref{sec:wire-conformance} derive the
exact counts for roots, nonces, proofs, rows, and query-dependent paths and
require logical and byte EOF at every recursive and outer boundary.&
Implement the v1 omission and exhaustion checks in both host and recursive
verification.  Until that implementation work lands, the current Rust path is
nonconforming; no alternate v1 field or suffix behavior is permitted.\\
\bottomrule
\end{longtable}

\endgroup
